%%%%%%%%%%%%%%%%%%%%%%%%%%%%%%%%%%%%%%%%%%%%%%%%%
\section{Object Detection for Wildlife Monitoring}\label{sec:lit_object_detection}

Modern object detection descends from two architecturally distinct lineages that emerged within a few years of each other. The two-stage, region-proposal lineage begins with \textit{R-CNN} \cite{girshickRichFeatureHierarchies2014a}, which paired an external region-proposal algorithm, \textit{Selective Search} \cite{uijlingsSelectiveSearchObject2013}, with a convolutional network applied independently to each proposed region; \textit{Fast R-CNN} \cite{girshickFastRCNN2015a} shared the convolutional feature computation across all proposals within a single network pass, and \textit{Faster R-CNN} \cite{renFasterRCNNRealTime2017} replaced the external proposal algorithm entirely with a learned Region Proposal Network, making the whole pipeline trainable end-to-end. \textit{Feature Pyramid Networks} \cite{linFeaturePyramidNetworks2017} later added a top-down, multi-scale feature hierarchy to this family to address the difficulty of detecting objects at greatly different scales -- a difficulty directly relevant to wildlife imagery, where the same species can occupy anywhere from a few pixels to most of the frame depending on distance to the camera. In parallel, a single-stage lineage discarded the separate proposal stage altogether: \textit{YOLO} \cite{redmonYouOnlyLook2016a} reframed detection as a single regression problem over the whole image, and \textit{SSD} \cite{liuSSDSingleShot2016a} added multi-scale default-box predictions to narrow the resulting accuracy gap while retaining a single forward pass. \textit{Focal Loss} \cite{linFocalLossDense2017} subsequently closed most of the remaining accuracy gap between the two lineages by explicitly correcting for the foreground--background class imbalance inherent to single-stage, dense prediction. For a comprehensive account of this more than twenty-year evolution beyond what is summarized here, see \cite{zouObjectDetection202023}.

Two-stage detectors retain a modest accuracy advantage in some regimes at the cost of two sequential network evaluations per image; single-stage detectors trade a small amount of accuracy for a fixed, single-pass inference cost that does not scale with the number of candidate objects in a scene. Because the target hardware for this thesis (\Cref{sec:model_selection}) has a fixed, tight per-frame latency budget, the single-stage lineage is the only one considered further in this thesis; \Cref{sec:lit_yolo_family} traces its dominant branch, the YOLO family, from the YOLOv5 generation actually trained here forward to the current state of the family.

%%%%%%%%%%%%%%%%%%%%%%%%%%%%%%%%%%%%%%%%%%%%%%%%%
\subsection{Challenges in Wildlife Object Detection}\label{sec:lit_wildlife_challenges}

Beyond the general architectural trajectory above, wildlife species detection poses four recurring challenges that motivate specific choices made elsewhere in this thesis.

First, distinguishing wildlife species is fundamentally a fine-grained visual categorization problem: unlike COCO-style categories (car, dog, person), many pairs of species differ only in subtle, localized visual cues rather than in gross shape or context \cite{weiFineGrainedImageAnalysis2022}. This small inter-class variation compounds a second, related problem: several groups of species that this thesis's target taxonomy must distinguish are visually near-identical to a general-purpose detector -- the plains-zebra/Grévy's-zebra, Eurasian-lynx/bobcat, and small-gazelle groupings later formalized as confusion clusters in this thesis's own evaluation methodology (\Cref{sec:evaluation_framework}) are instances of exactly the intra-order visual homogeneity that fine-grained recognition research identifies as its core difficulty \cite{weiFineGrainedImageAnalysis2022}.

Third, ecological image collections are long-tailed by nature: a handful of common, easily-photographed species dominate any naturally-collected corpus, while most species are represented by comparatively few images \cite{cuiClassBalancedLossBased2019,zhangDeepLongTailedLearning2023} -- the same imbalance this thesis's own sourcing effort encountered species-by-species (\Cref{sec:corpus_composition}).

Fourth, and specific to the wildlife domain, models trained on one visual domain generalize poorly to another. \textit{Recognition in Terra Incognita} \cite{beeryRecognitionTerraIncognita2018} showed that camera-trap classifiers and detectors that perform well at their training locations degrade sharply at new, unseen camera locations, since background, lighting, and camera geometry are confounded with location rather than with the animal itself. Larger deployed camera-trap classifiers reinforce the same pattern: cross-project transfer between ecologically similar studies still costs measurable accuracy \cite{williIdentifyingAnimalSpecies2019}, and cross-continental transfer can cost considerably more -- one large-scale camera-trap classifier's species-level accuracy fell from $98\%$ in-region to $82\%$ on an out-of-sample Canadian ungulate dataset \cite{tabakMachineLearningClassify2019} -- despite the same underlying kind of network having demonstrated near-human accuracy on its own training distribution ($>93.8\%$ species-identification accuracy, rising to $96.6\%$ on the confidently-classified subset) \cite{norouzzadehAutomaticallyIdentifyingCounting2018}. This is precisely the domain gap that motivated rejecting \textit{LILA BC}'s camera-trap archives as a training-data source for this thesis (\Cref{sec:sourcing_execution}): camera traps are static, fixed-position, frequently infrared-triggered sensors, whereas the target deployment is a handheld, daylight, close-to-mid-range binocular observation scenario, and the two distributions differ in exactly the location-, lighting-, and viewpoint-confounded ways this line of work documents. More broadly, the wildlife-monitoring machine-learning community has converged on the view that such domain shift, rather than raw model capacity, is now a primary obstacle to deploying automated recognition at scale \cite{tuiaPerspectivesMachineLearning2022}.

%%%%%%%%%%%%%%%%%%%%%%%%%%%%%%%%%%%%%%%%%%%%%%%%%
\subsection{The YOLO Model Family: YOLOv5 to YOLO26}\label{sec:lit_yolo_family}

Within the single-stage lineage, the YOLO family evolved along two intertwined axes: from anchor-based to anchor-free box prediction, and from non-maximum-suppression (NMS)-dependent to increasingly NMS-free inference. \textit{YOLOv1} \cite{redmonYouOnlyLook2016a} predicted bounding boxes directly as a regression target with no anchors at all, but its coarse output grid limited localization accuracy; \textit{YOLOv2}/\textit{YOLO9000} \cite{redmonYOLO9000BetterFaster2017} reintroduced anchor boxes, this time learned via $k$-means clustering over the training-set box distribution rather than hand-designed, and additionally proposed a hierarchical \textit{WordTree} label space that lets a single detector combine detection- and classification-only training data across a taxonomy -- an early precedent for the confidence-gated, taxonomy-aware inference this thesis revisits in a different system later in this chapter (\Cref{sec:lit_md_speciesnet}). \textit{YOLOv3} \cite{redmonYOLOv3IncrementalImprovement2018} deepened the backbone to Darknet-53 and added three-scale prediction heads for better small-object recall. The subsequent YOLOv4/\textit{YOLOv5} generation consolidated a Cross-Stage-Partial-Network (CSPNet) backbone with a PANet neck, and it is specifically the \textit{YOLOv5} \cite{jocher_yolov5_2020} codebase -- at the commit predating its relicensing under non-permissive commercial terms (\Cref{sec:model_selection}) -- that this thesis trains throughout \Cref{sec:augmentation_setups}. For an independent comparative review of the YOLOv2--YOLOv7 generations specifically, see \cite{nazirYouOnlyLook2023}.

From YOLOv5 onward, the family's dominant trajectory shifted from anchor-based to anchor-free prediction and, later, away from NMS altogether. \textit{YOLOX} \cite{geYOLOXExceedingYOLO2021} was the first widely-adopted YOLO variant to drop anchor boxes entirely, pairing a decoupled classification/regression head with the SimOTA dynamic label-assignment strategy; at matched, very low parameter counts (its YOLO-Nano configuration reaches $25.3\%$ COCO AP at $0.91\,\mathrm{M}$ parameters) it also offers one of the few controlled anchor-free comparisons available in this size regime. \textit{YOLOv6} \cite{liYOLOv6SingleStageObject2022} carried the anchor-free, industrial-deployment focus further with a re-parameterizable backbone aimed explicitly at throughput on commodity hardware. \textit{YOLOv9} \cite{wangYOLOv9LearningWhat2025} addressed information loss across network depth with Programmable Gradient Information and a Generalized Efficient Layer Aggregation Network (GELAN), reporting that a from-scratch GELAN model can match detectors pretrained on much larger datasets. \textit{YOLOv10} \cite{wangYOLOv10RealTimeEndtoEnd2024} completed the shift away from post-processing dependence by introducing consistent dual label assignments during training that make NMS unnecessary at inference time, removing a latency-variable, sequential step that had persisted through every earlier generation. \textit{YOLOv11} \cite{khanamYOLOv11OverviewKey2024} replaced YOLOv8's C2f block with the more compute-efficient C3k2 block and added a C2PSA spatial-attention module that YOLOv8 lacks, while \textit{YOLOv12} \cite{tianYOLOv12AttentionCentricRealTime2025} went further still, building an attention-centric architecture that matches the inference speed of the preceding CNN-based releases while outperforming YOLOv10-N and YOLOv11-N in accuracy at comparable latency. As of this writing, \textit{YOLO26} (Ultralytics, 2026) is the most recent release in the family; no independent architectural analysis or peer-reviewed description of YOLO26 comparable to those available for YOLOv9 through YOLOv12 has yet appeared in the surveyed literature, so this thesis makes no specific technical claim about it beyond noting that it continues the same anchor-free, NMS-free, export-oriented trajectory already established from YOLOv10 onward. It is worth stating plainly that none of the family's releases from YOLOv5 through YOLO26 -- including YOLOv8, YOLOv11, and YOLO26 itself -- has a peer-reviewed paper of its own; where cited above, the citation is to the codebase/release note or, where available, an independent architectural-analysis paper, which is standard practice for this literature. This thesis remains fixed on the YOLOv5 generation not for lack of awareness of these later releases, but for the commercial-licensing reason given in \Cref{sec:model_selection}.

The nano-scale variants most relevant to embedded deployment -- YOLOv8n, YOLOv10n, and YOLOv11n -- sit in the $2.6$--$3.2\,\mathrm{M}$-parameter range (\Cref{sec:target_class_definition}), and it is against this parameter budget, not against COCO's fixed $80$-class label space, that the 225-class taxonomy budget (\Cref{sec:target_class_definition}) was ultimately set. It should be stated plainly that no paper in the surveyed literature directly establishes a class-count-vs-accuracy ceiling for nano-scale detector backbones. The closest indirect evidence is capacity-accuracy comparisons at a fixed, comparatively small class count: \textit{PicoDet} shows a $0.99\,\mathrm{M}$-parameter detector reaching $30.6\%$ COCO AP, an absolute $4.8$-point improvement over the $0.91\,\mathrm{M}$-parameter YOLOX-Nano configuration at $25.3\%$ AP \cite{yuPPPicoDetBetterRealTime2021,geYOLOXExceedingYOLO2021}. This is evidence that discriminative capacity is a binding constraint for sub-million-parameter backbones in general, but it is a comparison across backbone size at one fixed class count, not a study of how that capacity trades off against a growing number of classes. The $200$--$250$-class ceiling adopted in \Cref{sec:target_class_definition} is therefore this thesis's own extrapolation from that general capacity-accuracy relationship, not a literature-established figure, and it is treated as such throughout.

%%%%%%%%%%%%%%%%%%%%%%%%%%%%%%%%%%%%%%%%%%%%%%%%%
\subsection{The MegaDetector + SpeciesNet Ensemble}\label{sec:lit_md_speciesnet}

A second architecture, already used elsewhere in this thesis's own pipeline as a pseudo-labeling and quality-filtering tool rather than as a deployment target (\Cref{sec:staged_filtering_pipeline,sec:speciesnet_matching}), takes the opposite structural approach from a single end-to-end detector: a species-agnostic localizer feeding a separate, much larger species classifier. \textit{MegaDetector} \cite{beery_megadetector_2019} is a general-purpose detector trained to localize animals, humans, and vehicles in camera-style imagery without attempting to resolve species identity itself; its output crops are then passed to \textit{SpeciesNet} \cite{gadot_speciesnet_2024}, a classifier built on an \textit{EfficientNetV2-M} backbone and trained on tens of millions of camera-trap images (on the order of $36\,\mathrm{M}$ training images spanning more than $2{,}000$ taxonomic labels in the largest reported training configuration) \cite{gadotCropNotCrop2024}. Rather than committing to a single species-level prediction regardless of confidence, this classification stage applies a confidence-gated taxonomic rollup: when the top species-level prediction falls below a confidence threshold, per-taxon confidence scores are summed upward across the taxonomy until some higher clade (genus, family, \ldots) clears the threshold, so that a low-confidence individual prediction degrades gracefully into a still-informative coarser one rather than an outright wrong species call \cite{gadotCropNotCrop2024}. A complementary geofencing step further constrains predictions to taxa expected in the image's region, rolling up to a coarser clade whenever the top prediction falls outside a region-specific allow-list \cite{gadotCropNotCrop2024}.

This ensemble represents something close to a practical accuracy ceiling for wildlife species recognition from camera-style imagery. Incorporating a detection stage ahead of classification, and cropping to the detected animal before classifying, was shown to yield a macro-averaged F1 improvement of roughly $25\%$ over whole-image classification on a large, long-tailed camera-trap dataset, a result that reproduced across a second public dataset and a smaller public benchmark \cite{gadotCropNotCrop2024}. This accuracy, however, comes at a combined computational cost that the two components' separate design does not hide: a full object detector and a separate, large classification backbone must both run, sequentially, once per detected crop in the image, rather than a class prediction falling out of a single localization pass. On general-purpose server hardware this cost is usually accepted given the accuracy gain, but it is difficult to reconcile with the QCS605's Hexagon 685 DSP budget: an \textit{EfficientNetV2-M} classifier alone is substantially larger than the nano-scale backbones ($1$--$3\,\mathrm{M}$ parameters, \Cref{sec:target_class_definition}) already established as computationally demanding to fit within this thesis's per-frame latency budget, before counting MegaDetector's own detector backbone or the per-crop inference multiplier once more than one animal is present in a frame. Deploying this ensemble directly on the target hardware would therefore require a substantial compression step -- for example, distilling SpeciesNet's classification knowledge into a nano-scale student, a design space explored at foundational depth in \Cref{sec:lit_kd_and_embedded} -- rather than being deployable as-is.

%%%%%%%%%%%%%%%%%%%%%%%%%%%%%%%%%%%%%%%%%%%%%%%%%
\subsection{End-to-End Detection vs.\ the Two-Stage Ensemble}\label{sec:lit_yolo_vs_ensemble}

The two preceding subsections set up a direct architectural choice: train a single, end-to-end YOLO-family detector directly on the 225-class taxonomy (\Cref{sec:lit_yolo_family}), or instead reuse or compress the \textit{MegaDetector}+\textit{SpeciesNet} ensemble (\Cref{sec:lit_md_speciesnet}) for on-device deployment. Each side of this choice trades a specific advantage against a specific cost.

The ensemble's principal advantage is higher accuracy together with graceful degradation. Cropping to a detected animal before classification measurably improves species-level accuracy over whole-image classification \cite{gadotCropNotCrop2024}, and the ensemble's confidence-gated taxonomic rollup means a classification that is not confident enough to commit to a species still yields a correct, if coarser, answer at the genus or family level rather than an outright wrong one \cite{gadotCropNotCrop2024}. A single end-to-end detector trained directly on species-level classes has no equivalent fallback built in by default: below its confidence threshold, a detection is simply dropped or misclassified, with no intermediate taxonomic level to fall back to unless one is deliberately engineered into the label space or loss function -- YOLO9000's WordTree hierarchy demonstrated that this is possible in principle for a detector, though it was not adopted for the 225-class taxonomy trained in this thesis \cite{redmonYOLO9000BetterFaster2017}.

The end-to-end detector's principal advantage is exactly the opposite side of the ensemble's principal cost. A single-stage detector performs one forward pass per frame at a fixed computational cost, independent of how many animals the frame contains, and exports through a single, well-understood embedded pipeline (\Cref{sec:lit_embedded_inference}). The \textit{MegaDetector}+\textit{SpeciesNet} ensemble, by contrast, pays a two-stage latency cost: \textit{MegaDetector} must localize before any classification can begin, and the classifier then runs once per detected crop, so total inference time scales with the number of detected instances rather than remaining constant per frame -- a cost structure that is particularly unfavourable under a fixed real-time budget such as the one this thesis targets on the QCS605. Compressing the ensemble to fit that budget would additionally require compressing \emph{two} separately-trained networks (detector and classifier) rather than one, each with its own accuracy-vs-size trade-off curve, whereas a single detector trained directly on the target taxonomy has only one network to compress.

This comparison is the literature grounding for the choice, stated concretely in \Cref{sec:model_selection}, to train a single end-to-end YOLO-family detector on the 225-class taxonomy rather than attempting to compress or port the \textit{MegaDetector}+\textit{SpeciesNet} ensemble itself onto the target hardware.

%%%%%%%%%%%%%%%%%%%%%%%%%%%%%%%%%%%%%%%%%%%%%%%%%
\subsection{Model Selection: Practical Constraints}\label{sec:model_selection}

Unlike the rest of this chapter, this subsection is not a literature survey: no dedicated model-selection methodology was carried out, and the choice of \textit{YOLOv5s} as the architecture trained throughout this thesis (\Cref{sec:augmentation_setups}) instead follows directly from three practical constraints established above and elsewhere in this chapter and in the Methods and Implementation chapter.

First, licensing fixed which \textit{YOLOv5} codebase revision was usable at all. \textit{YOLOv5} \cite{jocher_yolov5_2020} is only commercially usable up to commit \texttt{5cdad89}; later commits in the Ultralytics repository require additional licensing, which is incompatible with this thesis's requirement that the resulting weights be usable in a commercial product. This constraint alone rules out training on any newer Ultralytics release under the same license terms, independent of any architectural comparison among the releases surveyed in \Cref{sec:lit_yolo_family}.

Second, the combination of the $200$--$250$-class ceiling discussed in \Cref{sec:lit_yolo_family} and \Cref{sec:target_class_definition} with the QCS605's Hexagon 685 DSP compute budget rules out the \textit{MegaDetector}+\textit{SpeciesNet} ensemble (\Cref{sec:lit_yolo_vs_ensemble}) as a deployment target without a compression step -- such as distilling the ensemble's classification knowledge into a nano-scale student -- that this thesis does not undertake.

Third, the time budget available for this thesis, reflected in the reduced knowledge-distillation and embedded-inference scope reported in \Cref{chapter:results}, ruled out a comparative empirical architecture search. Two candidate directions were considered early in the project and then set aside on time grounds rather than on measured evidence: newer releases within the YOLO family itself (YOLOv8n, YOLOv10n, YOLOv11n, and later; \Cref{sec:lit_yolo_family}), and non-YOLO nano-scale detectors surveyed during early project scoping -- \textit{PicoDet} \cite{yuPPPicoDetBetterRealTime2021}, \textit{NanoDet}/\textit{NanoDet-Plus} \cite{rangilyuRangiLyuNanodet2026}, and \textit{EfficientDet-Lite} -- documented in \texttt{docs/2026-03-10\_object-detection-models-for-embedded-systems.md}. None of these alternatives were benchmarked empirically against \textit{YOLOv5s} within this thesis.

\textit{YOLOv5s} was therefore a practical, constraint-satisfying choice -- commercially licensable, small enough for the target hardware, and already integrated into the available training tooling -- rather than the winner of an architecture comparison against newer YOLO releases or non-YOLO nano detectors. This is stated here plainly and treated as a named limitation of the work (\Cref{sec:limitations}), not a hidden one.
